%% report.tex — RV Forecasting Benchmark
%% Compile with: pdflatex report.tex  (twice for cross-references)
\documentclass[11pt, a4paper]{article}

\usepackage[top=2.2cm, bottom=2.2cm, left=2.5cm, right=2.5cm]{geometry}
\usepackage{amsmath, amssymb, bm}
\usepackage{booktabs, multirow, array}
\usepackage[ruled, linesnumbered, noend]{algorithm2e}
\usepackage{hyperref}
\usepackage{xcolor}
\usepackage{parskip}
\usepackage{microtype}

\hypersetup{colorlinks=true, linkcolor=blue!60!black, citecolor=blue!60!black}
\newcolumntype{C}[1]{>{\centering\arraybackslash}p{#1}}

%% notation shortcuts
\newcommand{\RV}{\mathrm{RV}}
\newcommand{\E}{\mathbb{E}}
\newcommand{\R}{\mathbb{R}}
\newcommand{\hath}{\hat{y}_{t+h}}
\newcommand{\xt}{x_t}
\newcommand{\st}{s_t}
\newcommand{\W}{w}

\title{\textbf{Realised Volatility Forecasting — Multi-Asset Benchmark}\\[4pt]
       \large Direct Multi-Step Models: HAR, NLinear, DLinear, SAS, RNN}
\author{}
\date{\today}

\begin{document}
\maketitle
\vspace{-1em}

% ─────────────────────────────────────────────────────────────────────────────
\section{Setup}
% ─────────────────────────────────────────────────────────────────────────────

\textbf{Input series.} Either raw $y_t = \RV_t > 0$ (default) or
log-transformed $y_t = \log\RV_t \in \mathbb{R}$ (\texttt{--log}).
All internal scaling (z-scoring) is handled inside each model class.

\textbf{Forecasting task.} Direct multi-step: one readout per
$h \in \{1,5,10\}$ days.  Readouts are ridge-regression with
$\alpha \in \{10^k : k \in [-4,5],\,\Delta k{=}0.25\}$ (37 values),
5-fold time-series CV.

% ─────────────────────────────────────────────────────────────────────────────
\section{Models}
% ─────────────────────────────────────────────────────────────────────────────

\subsection{HAR — Heterogeneous Autoregressive Model (Corsi 2009)}

\textbf{Features.} $\bar{y}^{(w)}_t = \tfrac{1}{5}\sum_{k=0}^{4}y_{t-k}$,\quad
$\bar{y}^{(m)}_t = \tfrac{1}{22}\sum_{k=0}^{21}y_{t-k}$.

\begin{equation}
  \hat{y}_{t+h} = \alpha + \beta_d\,y_t + \beta_w\,\bar{y}^{(w)}_t
                         + \beta_m\,\bar{y}^{(m)}_t
\end{equation}

\textbf{HAR-OLS}: $\min_\beta \|X\beta - y\|_2^2$.\quad
\textbf{HAR-Ridge}: column-normalise features; penalise slopes only;
back-transform coefficients after fit.

% ─────────────────────────────────────────────────────────────────────────────
\subsection{NLinear and DLinear (Zeng et al.\ 2023)}

Lookback window $\bm{y}_t = (y_{t-L+1},\ldots,y_t)^\top \in \R^L$.

\textbf{NLinear} — subtract last value (removes level, keeps increments):
\begin{equation}
  \tilde{\bm{y}}_t = \bm{y}_t - y_t, \qquad
  \hat{y}_{t+h} = \tilde{\bm{y}}_t^\top \W + y_t
\end{equation}

\textbf{DLinear} — trend/seasonality decomposition:
\begin{equation}
  \bm{\tau}_t = \mathrm{MA}_k(\bm{y}_t), \quad
  \bm{\varsigma}_t = \bm{y}_t - \bm{\tau}_t, \qquad
  \hat{y}_{t+h} = \begin{bmatrix}\bm{\tau}_t \\ \bm{\varsigma}_t\end{bmatrix}^\top \W
\end{equation}

% ─────────────────────────────────────────────────────────────────────────────
\subsection{GARCH — AR($p$)-GARCH(1,1)}

$y_t = \mu + \sum_{i=1}^{p} \phi_i y_{t-i} + \varepsilon_t$,\quad
$\sigma^2_t = \omega + \alpha_{\scriptscriptstyle G}\varepsilon^2_{t-1}
            + \beta_{\scriptscriptstyle G}\sigma^2_{t-1}$.
Multi-step: iterate variance recursion analytically. Estimated by QMLE.

% ─────────────────────────────────────────────────────────────────────────────
\subsection{SAS — State Affine System}

\textbf{Diagonal basis.}
$P_0,\ldots,P_p \in \R^n$,\; $Q_0,\ldots,Q_q \in \R^n$, input $z\in\mathbb R$.

\begin{align}
  P(z) &= \mathrm{clip}\!\Bigl(\sum_{d=0}^{p} P_d\,z^d,\;
           {-1{+}\epsilon},\;{1{-}\epsilon}\Bigr) \in (-1,1)^n \\
  Q(z) &= \sum_{d=0}^{q} Q_d\,z^d \in \R^n \\
  s_t  &= P(z_t) \odot s_{t-1} + Q(z_t)
\end{align}

Clipping guarantees $|P(z)_i|<1$ (unconditional stability).

\textbf{Standard readout.}
\begin{equation}
  \hat{y}_{t+h} = s_t^\top \W_h, \qquad
  \W_h = \operatorname*{arg\,min}_{w} \|S\,w - Y_h\|^2 + \alpha\|w\|^2
\end{equation}
where $S \in \R^{(T-p-h)\times n}$ are post-washout states and
$Y_h \in \R^{T-p-h}$ are z-scored targets $y_{t+h}$.

\textbf{Residual-target variant (SAS\textsubscript{res}).}
Instead of fitting the readout on the level $y_{t+h}$, fit it on the
increment $r_{t,h} = y_{t+h} - y_t$ (per-horizon z-scored separately):

\begin{equation}
  \W_h = \operatorname*{arg\,min}_{w}
         \bigl\|S\,w - \tilde{r}_h\bigr\|^2 + \alpha\|w\|^2,
  \qquad \tilde{r}_h = \frac{r_h - \mu_{r,h}}{\sigma_{r,h}}
\end{equation}

At prediction time reconstruct the level using the anchor $y_t$:

\begin{equation}
  \hat{r}_{t,h} = \sigma_{r,h}\,(s_t^\top \W_h) + \mu_{r,h},
  \qquad \hat{y}_{t+h} = y_t + \hat{r}_{t,h}
\end{equation}

The anchor $y_t > 0$ prevents near-zero predictions and stabilises QLIKE.
Scalers $(\mu_{r,h}, \sigma_{r,h})$ are computed per horizon on the training
window and updated at each refit.

\textbf{Initialisation of $P$.}
Let $\rho$ be the spectral-norm bound and $\delta = \max(0.005,\,(1{-}\rho)\cdot 0.3)$.
\begin{itemize}\itemsep1pt
  \item $P_0[i] \overset{\text{iid}}{\sim} \mathcal{U}(-\rho,\,\rho)$
  \item $P_k[i] \overset{\text{iid}}{\sim} \mathcal{N}(0,\,\delta/5^k)$
        for $k \geq 1$
\end{itemize}

\textbf{Initialisation of $Q$.}
To keep injection variance uniform across degrees (for $z \sim \mathcal{N}(0,1)$,
$\mathbb{E}[z^{2k}] = (2k-1)!!$):
\begin{equation}
  Q_k[i] \overset{\text{iid}}{\sim}
  \mathcal{N}\!\left(0,\;\frac{1}{\sqrt{n\,(2k-1)!!}}\right)
\end{equation}

\textbf{Associative structure.} Composition rule on pairs $\ell_t = (P(z_t),Q(z_t))$:
\begin{equation}
  (A_j,\,b_j) \circ (A_i,\,b_i) =
  \bigl(A_j \odot A_i,\;\;A_j \odot b_i + b_j\bigr)
\end{equation}

This is associative $\Rightarrow$ parallel prefix scan (depth $O(\log T)$
vs.\ $O(T)$ sequential).

\begin{algorithm}[H]
\caption{Two-Level Parallel Scan — compute $s_1,\ldots,s_T$}
\SetKwInOut{Input}{Input}\SetKwInOut{Output}{Output}
\Input{Transitions $\ell_1,\ldots,\ell_T$;\; chunk size $B$;\; $s_0 = \bm{0}$}
\Output{States $s_1,\ldots,s_T$}
$K \leftarrow \lceil T/B \rceil$; pad to $KB$ \\
\textbf{Phase 1} (parallel over $K$ chunks):
  $(A^\text{intra}_{k,\cdot},b^\text{intra}_{k,\cdot})
  \leftarrow \textsc{AssocScan}(\circ,\;\ell_{kB+1},\ldots,\ell_{(k+1)B})$ \\
\textbf{Phase 2} (inter-chunk scan over $K$ endpoints):
  $(A^\text{inter},b^\text{inter})
  \leftarrow \textsc{AssocScan}(\circ,\;(A^\text{intra}_{1,B},b^\text{intra}_{1,B}),\ldots)$ \\
\textbf{Phase 3} (carry):
  $c_0 = s_0$;\;
  $c_k = A^\text{inter}_{k-1} \odot s_0 + b^\text{inter}_{k-1}$ \\
\textbf{Phase 4} (parallel over $K$ chunks):
  $s_{kB+j} = A^\text{intra}_{k,j} \odot c_k + b^\text{intra}_{k,j}$
\end{algorithm}

% ─────────────────────────────────────────────────────────────────────────────
\subsection{RNN — LSTM / GRU / LRU}

Readout: $\hat{y}_{t+h} = \text{Linear}_h(h_T)$, one head per $h$.
Training: Adam, early stopping (patience = 20, val = 20\% chronological).

\textbf{LRU} (Orvieto et al.\ 2023):
\begin{equation}
  h_t = \sigma(\nu) \odot h_{t-1} + Bz_t, \qquad y_t = Ch_t,
  \qquad \sigma(\nu) \in (0,1)^d
\end{equation}

% ─────────────────────────────────────────────────────────────────────────────
\section{Evaluation Protocol}
% ─────────────────────────────────────────────────────────────────────────────

\subsection{Rolling Out-of-Sample Loop}

Each model owns its preprocessing (z-scoring, log-transform).
The loop passes raw series values unchanged.

\begin{algorithm}[H]
\caption{Rolling OOS Evaluation}
\SetKwInOut{Input}{Input}\SetKwInOut{Output}{Output}
\Input{Series $y_1,\ldots,y_T$; window $W$; refit freq $R$; horizons $\mathcal{H}$}
\Output{Records $\{(\hat{y}_{t+h,m},\,y_{t+h},\,\mathrm{sqerr},\,\mathrm{qlike})\}$}
\For{$t = W$ \KwTo $T - \max(\mathcal{H})$}{
  \If{steps since last refit $\geq R$}{
    \ForEach{model $m$}{$m.\textsc{Fit}(y_{t-W:t},\;\mathcal{H})$} \tcp*{model z-scores internally}
  }
  \ForEach{model $m$}{$m.\textsc{Update}(y_t)$} \tcp*{advance state to $t$}
  \ForEach{model $m$, $h \in \mathcal{H}$}{
    $\hat{y}_{t+h} \leftarrow m.\textsc{Predict}(h)$ \tcp*{in original scale}
    record $\bigl(\hat{y}_{t+h},\;y_{t+h},\;
      (\hat{y}_{t+h}-y_{t+h})^2,\;
      \mathcal{L}(\hat{y}_{t+h},y_{t+h})\bigr)$\;
  }
}
\end{algorithm}

\subsection{Loss Functions}

\textbf{Mean Squared Error.}
\begin{equation}
  \mathrm{MSE}^{(h)}_m = \frac{1}{T_{\mathrm{oos}}}
    \sum_{t=W}^{T-h} \bigl(\hat{y}_{t+h,m} - y_{t+h}\bigr)^2
\end{equation}

\textbf{QLIKE} (Patton 2011) — level-space ($y_t > 0$, default):
\begin{equation}
  \mathcal{L}(y,\hat{y}) = \frac{y}{\hat{y}} - \log\frac{y}{\hat{y}} - 1
  \qquad (\hat{y} \text{ clipped to } 10^{-8})
\end{equation}
$\mathcal{L} \geq 0$ with minimum 0 at $\hat{y}=y$.
Scale-free: $\mathcal{L}(ky,k\hat{y}) = \mathcal{L}(y,\hat{y})$.
Strongly penalises underforecasting ($\hat{y} \ll y$).

\textbf{QLIKE — log-space} (\texttt{--log} mode, $y_t \in \mathbb{R}$):
Apply the level formula to $e^{y_t},\,e^{\hat{y}_t}$:
\begin{equation}
  \mathcal{L}_{\log}(y,\hat{y})
    = \frac{e^{\hat{y}}}{e^{y}} - \left(y - \hat{y}\right) - 1
    = e^{\hat{y}-y} + (\hat{y}-y) - 1
\end{equation}
Valid for any finite $y,\hat{y}$ (no positivity constraint required).

\subsection{Model Confidence Set (Hansen et al.\ 2011)}

Given loss matrix $L \in \R^{T\times M}$ (raw squared errors, one column per model):
\begin{enumerate}\itemsep0pt
  \item $d_{t,ij} = L_{t,i} - L_{t,j}$;\quad
        $\bar{d}_{ij} = T^{-1}\!\sum_t d_{t,ij}$
  \item $T_R = \max_{m\neq b}\; \bar{d}_{mb}/\hat{\sigma}_{mb}$
        \quad ($\hat{\sigma}^2_{mb}$: circular block bootstrap, $B$ reps,
        block width $w = \lfloor T^{1/3} \rfloor$)
  \item Eliminate the worst model if $p\text{-value}(T_R) < \alpha$; repeat.
  \item $\mathrm{MCS}_\alpha$ = surviving set.
\end{enumerate}

\textbf{MCS Frequency.} Run MCS independently per symbol per horizon;
report survival rate $c_m^{(h)} / N_{\mathrm{sym}}$.

\subsection{Diebold-Mariano Test (Harvey-Leybourne-Newbold 1997)}

For each (model, horizon, symbol), test whether a challenger beats the benchmark.
\textbf{Raw loss differences} are used for both MSE and QLIKE:
\begin{equation}
  d_t = \mathcal{L}_{t,\mathrm{challenger}} - \mathcal{L}_{t,\mathrm{benchmark}}
\end{equation}
$H_0\!: \E[d_t]=0$;\; $H_1^{<}\!: \E[d_t]<0$ (challenger better).

\textbf{HAC variance} with Newey-West, $h-1$ lags
(MA($h-1$) structure of $h$-step errors):
\begin{equation}
  \hat{V} = \hat{\gamma}_0 + 2\sum_{k=1}^{h-1}
            \left(1 - \frac{k}{h}\right)\hat{\gamma}_k,
  \qquad \hat{\gamma}_k = T^{-1}\sum_t (d_t-\bar d)(d_{t-k}-\bar d)
\end{equation}

\textbf{HLN finite-sample correction}:
\begin{equation}
  \mathrm{DM}^* = \sqrt{\frac{T+1-2h+h(h-1)/T}{T}}\;\cdot\;
                  \frac{\bar{d}}{\sqrt{\hat{V}/T}}
  \sim t_{T-1}
\end{equation}

Report: fraction of symbols where $p^* < \alpha$ (one-sided), per horizon.

% ─────────────────────────────────────────────────────────────────────────────
\section{Hyperparameters}
% ─────────────────────────────────────────────────────────────────────────────

\begin{table}[ht]
\centering
\caption{Fixed hyperparameters used in all experiments.}
\label{tab:hyperparam}
\begin{tabular}{llr}
\toprule
\textbf{Component} & \textbf{Parameter} & \textbf{Value} \\
\midrule
\multirow{3}{*}{Evaluation}
  & Training window $W$    & 2\,000 \\
  & Refit frequency $R$    & 20 steps ($\approx$1 month) \\
  & Horizons $\mathcal{H}$ & $\{1, 5, 10\}$ days \\
\midrule
\multirow{2}{*}{MCS}
  & Significance $\alpha$         & 0.10 \\
  & Bootstrap replications $B$    & 1\,000 \\
\midrule
\multirow{2}{*}{HAR}
  & Lags $(d, w, m)$ & $(1, 5, 22)$ days \\
  & Ridge CV folds   & 5 \\
\midrule
\multirow{2}{*}{NLinear / DLinear}
  & Lookback $L$      & 20 \\
  & DLinear MA kernel & 5 \\
\midrule
\multirow{2}{*}{GARCH}
  & AR lags $p$         & 5 \\
  & Estimation          & QMLE \\
\midrule
\multirow{6}{*}{SAS (diagonal)}
  & Reservoir size $n$            & 200 \\
  & Spectral norm $\rho$          & 0.95 \\
  & $P$-polynomial degree $p$     & 1 \\
  & $Q$-polynomial degree $q$     & 1 (SAS\textsubscript{q1}) or 2 (SAS\textsubscript{q2}) \\
  & Washout                       & 50 steps \\
  & Residual-target               & off (SAS) / on (SAS\textsubscript{res}) \\
\midrule
\multirow{5}{*}{RNN (GRU)}
  & Hidden size     & 32 \\
  & Layers          & 1 \\
  & Lookback        & 20 \\
  & Epochs (max)    & 500 \\
  & Early stopping  & patience = 20, val = 20\% \\
\midrule
\multirow{3}{*}{Ridge grid}
  & Range    & $10^k$, $k \in [-4, 5]$ \\
  & Spacing  & $\Delta k = 0.25$,\; 37 values \\
  & CV folds & 5 \\
\bottomrule
\end{tabular}
\end{table}

\newpage
\section{Results}


\end{document}
